\newpage
\phantomsection
\label{prf:congruence-modulo-integer-equivalence}
\LRAProofFor{thm:congruence-modulo-integer-equivalence}

\begin{remark*}[Return]
\hyperref[thm:congruence-modulo-integer-equivalence]{Return to Theorem}
\end{remark*}

\begin{theorem*}[Congruence Modulo an Integer Is an Equivalence Relation]
For each positive integer \(n\), congruence modulo \(n\) is an equivalence
relation on \(\mathbb{Z}\).
\end{theorem*}

\begin{proof}[Professional Standard Proof]
\LRAProofBodyStart
\textbf{Professional Standard Proof.}
Fix positive \(n\). Reflexivity holds because \(a-a=0=n\cdot 0\). Symmetry
holds because if \(a-b=nk\), then \(b-a=n(-k)\). Transitivity holds because
if \(a-b=nk\) and \(b-c=n\ell\), then \(a-c=(a-b)+(b-c)=n(k+\ell)\). Thus
congruence modulo \(n\) is reflexive, symmetric, and transitive.
\end{proof}

\begin{proof}[Detailed Learning Proof]
\LRAProofBodyStart
\textbf{Detailed Learning Proof.}
Let \(n\) be positive. For every integer \(a\), the difference \(a-a\) is
\(0\), and \(0\) is divisible by \(n\), so \(a\equiv a\pmod n\). If
\(a\equiv b\pmod n\), then \(a-b=nk\) for some integer \(k\). Multiplying by
\(-1\) gives \(b-a=n(-k)\), so \(b\equiv a\pmod n\). Finally, if
\(a\equiv b\pmod n\) and \(b\equiv c\pmod n\), write
\(a-b=nk\) and \(b-c=n\ell\). Adding gives
\(a-c=n(k+\ell)\), so \(a\equiv c\pmod n\). These are precisely the three
requirements for an equivalence relation.
\end{proof}

\begin{remark*}[Proof structure]
The three equivalence-relation laws are divisibility calculations in
\(\mathbb{Z}\).
\end{remark*}

\begin{dependencies}
\begin{itemize}
  \item \hyperref[def:equivalence-relation]{Equivalence Relation}
  \item \hyperref[def:congruence-modulo-integer]{Congruence Modulo an Integer}
  \item \hyperref[def:divisibility-in-z]{Divisibility in \(\mathbb{Z}\)}
\end{itemize}
\end{dependencies}

\clearpage
